% Partie combinaison et conclusion.
% Sources : results/comparison_local.csv, enquete/results/juliet/{juliet2_test,drs_recalib_test}.csv, enquete/combinaison.md,
% PLAN.md (score plateforme : 21,59 / 0,702 / 0,585 ; FID test du modèle rendu 21,15).
\section{Combining the levers, and conclusion}\label{sec:combination}

\textbf{Gains do not simply add up.} Each lever's effect depends on the model it is applied to. The new loss only pays off
once the generator is averaged (Sec.~\ref{sec:training}). Filtering with the original $D$ helps the baseline, with or
without EMA (30.0 to 24.8 with EMA), but not the final model, unless $D$ is recalibrated (Sec.~\ref{sec:inference}). The
data levers brought nothing to stack.

\textbf{Submitted model.} RpGAN + R1 ($\gamma=1$) with EMA, the seed with the best validation score, and no filter. It
scores 21.15 on our test set (the method averages $22.1\pm1.3$ over seeds) and, on the platform, FID 21.59 with
precision 0.70 and recall 0.59, against 38.7, 0.47 and 0.51 for the baseline. Its samples are shown in
Fig.~\ref{fig:samples}c. With the recalibrated OBRS filter, found after submission, we would expect about 2 points more
for three times more generated images.

\textbf{What we learned.}
\begin{itemize}
\item On a fixed, small generator, \emph{how} it is trained matters more than \emph{what} it sees: the training lever
gave all of the gain.
\item Ideas that work on large, diverse datasets do not automatically transfer to MNIST: selecting typical images
removes useful diversity, and noise does not cure a numerical problem.
\item A tool built on the discriminator (here, a filter) must be checked against the discriminator's actual ability:
the better the GAN, the more confused its $D$.
\item The choice of the measure shapes the conclusions: a digit classifier and Inception do not see the same defects,
so we always report which one we use, and look at the images.
\end{itemize}

\textbf{Limits.} Three seeds per variant only detect differences of about 3.6 points; some experiments ran on two
different machines (about 2 points apart); the data experiments used their own protocol; and two controls are missing
(R1 with the baseline loss, selection of a random half of the images).

\takeaway{Averaging the generator and a relativistic loss with R1 take the FID from 38.7 to 22.1 (21.59 on the
platform); filtering with a recalibrated discriminator could take it to about 19.7.}
